\input{libraries/report-preamble.tex}
\title{A Full-Range Binary64 Exponential}
\begin{document}
\maketitle

\section{Task and verification status}
\texttt{LeanExe.Lib.Transcendental.exp} computes $e^x$ for binary64 inputs.
The function accepts and returns IEEE 754 bit patterns in \texttt{UInt64}
words.  LeanExe compiles the Lean implementation and its clients to
WebAssembly (WASM), which the examples execute in Wasmtime.

The implementation ports Arm's 128-entry, degree-five exponential
algorithm~\cite{arm,data}.  It selects round-to-nearest, ties-to-even
arithmetic and separate multiplication and addition.  Arm reports a
worst-case error of $0.511$ ulp for this configuration without fused
multiply-add.  The port's target is an error below one ulp.

The theorem \texttt{Project.ExpArm.Spec.exp\_binary} proves that the recorded
WASM bytes decode and validate to a module whose exponential terminates for
every input word, preserves the complete store, and returns the result of an
exact binary64 computation of the algorithm.  The proof uses Talos's WASM
semantics with round-to-nearest, ties-to-even arithmetic and proves that
instantiation initializes the lookup table.  The theorem \texttt{exp\_exact}
applies to any store containing that table, supporting repeated calls.

For $|x|<2^{-54}$, \texttt{exp\_tiny\_error} proves that the WASM returns one
with absolute error below $2^{-53}$ against the real exponential.
The full-range numerical error proof remains outstanding.
Development tests compare WASM, the Lean port,
Lean's \texttt{Float.exp}, JavaScript's \texttt{Math.exp}, and a high-precision
decimal reference.  The largest measured error for the port on 49,821 inputs
is $0.504059$ ulp.

\section{Range reduction and evaluation}
The reduction chooses an integer $k$ near $128x/\log 2$ and writes
\[
  x=k\log(2)/128+r,\qquad e^x=2^{k/128}e^r.
\]
The implementation obtains $k$ through a binary64 addition and subtraction
of $1.5\cdot2^{52}$.  Splitting $-\log(2)/128$ into two constants reduces
the rounding error in $r$.  The reduced argument has magnitude near or
below $\log(2)/256$.

For $j=k\bmod128$, a table supplies a leading value $H_j$ and a correction
$t_j$ such that $2^{j/128}\approx H_j(1+t_j)$.  Its second word incorporates
an exponent adjustment, allowing the scale to be formed by integer bit
operations.  The table has 128 pairs of 64-bit words, totaling 2 KiB.

The polynomial evaluates the correction
\[
 t_j+r+r^2(C_2+rC_3)+r^4(C_4+rC_5).
\]
The coefficient words and evaluation order follow Arm's source.  Each
arithmetic operation rounds separately.  Normal reconstruction computes
\texttt{scale + scale * tmp}.

For large positive inputs, reconstruction first reduces the scale's exponent
by 1009, evaluates the correction, and multiplies by $2^{1009}$.  For large
negative inputs, it increases the scale's exponent by 1022 and multiplies
the corrected result by $2^{-1022}$.  When the corrected value is below one,
a compensated sum rounds it to the precision required by the subnormal
result before that multiplication.  This step controls double rounding.

\section{Input behavior and storage}
Both signed zeros return one.  Inputs with magnitude below $2^{-54}$ also
return one.  Positive infinity returns positive infinity.  Negative infinity
returns positive zero.  All NaNs return the canonical quiet NaN with word
\texttt{0x7FF8000000000000}.  Finite results can overflow to positive infinity
or underflow through subnormal values to positive zero.

The API specifies result values under round-to-nearest arithmetic.  It has
no error-number or floating-point exception-flag interface.  Native Lean
comparisons assume that the host uses round-to-nearest arithmetic and
preserves subnormals.

The source uses a named \texttt{Array UInt64} literal.  LeanExe places its
length and words in a WASM data segment and emits bounds-checked loads for
direct indexed reads.  The allocation arena begins after the static data,
and reset preserves the table.  The execution theorem proves that an
exponential call preserves all memory and globals.  Array updates retain the compiler's
ordinary ownership and allocation behavior.

\section{Executable implementation}
The listing comes from the component's \texttt{Basic.lean}.  The table is in
the adjacent \texttt{Table.lean}.  The helper \texttt{ExpArm.rescale} implements
the overflow and subnormal reconstruction described above.

\lstinputlisting[firstline=25,lastline=48]{LeanExe/Lib/Transcendental/ExpArm/Basic.lean}

\newpage
\section{Commands and a sigmoid client}
The commands run from the repository root after following the
\href{https://github.com/jsmorph/leanexe/blob/lib1/DEVELOPING.md}{setup guide}.

\begin{lstlisting}
tools/seminum exp 1
2.718281828459045
tools/seminum sigmoid 0
0.5
\end{lstlisting}

The decimal command accepts finite binary64 inputs.  The hexadecimal command
accepts every input word.  The sigmoid client computes
$\sigma(x)=1/(1+e^{-x})$ using $z=e^{-|x|}$: its result is $z/(1+z)$ for
negative inputs and $1/(1+z)$ otherwise.  This form preserves small positive
results for large negative inputs without an overflowing intermediate
exponential.  Both clients import and call the library function.

\section{Numerical tests}
\texttt{node test/exp.js} runs 49,821 cases.  It writes measurements to
\path{build/seminum/exp-accuracy.json}.  The corpus includes signed zeros,
infinities, quiet and signaling NaNs, random binary64 words, inputs distributed
across the finite output range, and neighboring words around reduction,
overflow, and underflow boundaries.  WASM and the Lean port agree bit for bit
on every case.  Separate tests cover the sigmoid client, command-line output,
and allocation counts.

The reference uses Python's standard-library \texttt{Decimal.exp}, which
rounds to nearest at the requested decimal precision~\cite{decimal}.
Conversion from the binary64 input is exact.  Calculation starts at 100
significant digits and increases precision until the neighboring decimal
bounds round to the same binary64 word.  The error calculation uses those
bounds.  One ulp is the binary64 spacing in the exact result's binade, with
spacing $2^{-1074}$ for subnormal results.

\begin{center}
\begin{tabular}{lrr}
\hline
Computation & Maximum error (ulp) & Differently rounded cases\\
\hline
LeanExe WASM & 0.504059 & 49\\
Lean \texttt{Float.exp} & 0.503354 & 36\\
JavaScript \texttt{Math.exp} & 0.871814 & 4153\\
\hline
\end{tabular}
\end{center}
The last column counts differences from the reference's rounded binary64
result, excluding NaN payloads.  Lean 4.34.0-rc2 implements
\texttt{Float.exp} through the external C function \texttt{exp}~\cite{lean}.
These measurements used Linux AArch64 and Node.js 24.13.0.  Runtime-library
comparisons depend on the installed implementations.  Throughput and latency
have yet to be measured.

\begin{thebibliography}{9}
\bibitem{arm} Arm,
\href{https://github.com/ARM-software/optimized-routines/blob/master/math/exp.c}{\emph{Double-precision exponential}}.
Source of the reduction, evaluation order, and reconstruction.  The port
retains Arm's copyright notice and uses the MIT license.
\bibitem{data} Arm,
\href{https://github.com/ARM-software/optimized-routines/blob/master/math/exp_data.c}{\emph{Shared exponential data}}.
Source of the table, split logarithm, coefficients, and reported error analysis.
\bibitem{lean} Lean 4.34.0-rc2,
\href{https://github.com/leanprover/lean4/blob/v4.34.0-rc2/src/Init/Data/Float/Float.lean}{\emph{Floating-point API}}.
Documents the external C implementation of \texttt{Float.exp}.
\bibitem{decimal} Python,
\href{https://docs.python.org/3/library/decimal.html#decimal.Decimal.exp}{\emph{Decimal arithmetic}}.
Documents exact conversion from binary floating point and correct rounding
of the exponential at the selected decimal precision.
\end{thebibliography}
\end{document}
